% ============================================================
%  PEMBAHASAN KUIS 3 MATEMATIKA (MAE)  T.P. 2026/2027
%  Tema: "National Card Strategy Championship (NCSC)"
%  Materi inti: Bab 6 Statistika (6.1-6.5), Bab 7 Peluang (7.1-7.5)
%  Data sinkron: rekap skor 40 peserta (x=25), set remi standar 52 kartu.
%  Kompilasi: pdfLaTeX (gambar berada di folder yang sama)
% ============================================================
\documentclass[11pt,a4paper]{article}

% ---------- PAKET ----------
\usepackage[utf8]{inputenc}
\usepackage[T1]{fontenc}
\usepackage[bahasa]{babel}
\usepackage{amsmath,amssymb,mathtools}
\usepackage{geometry}
\usepackage{xcolor}
\usepackage{graphicx}
\usepackage{enumitem}
\usepackage{tikz}
\usetikzlibrary{calc,arrows.meta}
\usepackage[most]{tcolorbox}
\usepackage{fancyhdr}
\usepackage{array,booktabs}
\usepackage{pifont}
\usepackage{icomma}
\usepackage[colorlinks=true,linkcolor=biru,urlcolor=biru]{hyperref}

\geometry{a4paper,margin=2.5cm,headheight=15pt}

% ---------- WARNA ----------
\definecolor{biru}{HTML}{1F4E79}
\definecolor{birumuda}{HTML}{D6E4F0}
\definecolor{hijau}{HTML}{2E7D32}
\definecolor{hijaumuda}{HTML}{E3F2E1}
\definecolor{oranye}{HTML}{C75B12}
\definecolor{oranyemuda}{HTML}{FBE8DA}
\definecolor{abu}{HTML}{F2F2F2}
\definecolor{ungu}{HTML}{6A1B9A}
\definecolor{ungumuda}{HTML}{EDE2F3}
\definecolor{merah}{HTML}{C62828}
\definecolor{merahmuda}{HTML}{FCE4E4}
\definecolor{tosca}{HTML}{00838F}
\definecolor{toscamuda}{HTML}{E0F2F1}

% ---------- HEADER / FOOTER ----------
\pagestyle{fancy}\fancyhf{}
\renewcommand{\headrulewidth}{0.8pt}
\fancyhead[L]{\small\textcolor{ungu}{Pembahasan Kuis 3: MAE}}
\fancyhead[R]{\small\textcolor{ungu}{T.P. 2026/2027}}
\fancyfoot[C]{\thepage}

% ---------- KOTAK: SOAL (ungu) ----------
\newtcolorbox{soalbox}[1]{
  enhanced, breakable, colback=ungumuda, colframe=ungu, boxrule=0.8pt,
  arc=3pt, left=10pt, right=10pt, top=8pt, bottom=8pt,
  title={\textbf{#1}}, fonttitle=\bfseries\color{white}, coltitle=white,
  attach boxed title to top left={xshift=10pt,yshift=-10pt},
  boxed title style={colback=ungu,arc=2pt}}

% ---------- KOTAK: KONSEP KUNCI (hijau) ----------
\newtcolorbox{ingat}[1]{
  enhanced, breakable, colback=hijaumuda, colframe=hijau, boxrule=0.6pt,
  arc=3pt, left=10pt, right=10pt, top=6pt, bottom=6pt,
  title={\textbf{Konsep Kunci: #1}}, fonttitle=\bfseries\color{white}, coltitle=white,
  attach boxed title to top left={xshift=10pt,yshift=-10pt},
  boxed title style={colback=hijau,arc=2pt}}

% ---------- KOTAK: JEBAKAN (merah) ----------
\newtcolorbox{jebakan}{
  enhanced, breakable, colback=merahmuda, colframe=merah, boxrule=0.6pt,
  arc=3pt, left=10pt, right=10pt, top=6pt, bottom=6pt,
  title={\textbf{\ding{55}\ Jebakan yang Sering Terjadi}},
  fonttitle=\bfseries\color{white}, coltitle=white,
  attach boxed title to top left={xshift=10pt,yshift=-10pt},
  boxed title style={colback=merah,arc=2pt}}

% ---------- KOTAK: TIPS (tosca) ----------
\newtcolorbox{tips}{
  enhanced, breakable, colback=toscamuda, colframe=tosca, boxrule=0.6pt,
  arc=3pt, left=10pt, right=10pt, top=6pt, bottom=6pt,
  borderline west={3pt}{0pt}{tosca}}

% ---------- PERINTAH BANTU ----------
\newcommand{\jawab}{\par\smallskip\noindent\textit{\textcolor{oranye}{Penyelesaian:}}\par\smallskip}
\newcommand{\langkah}[1]{\par\smallskip\noindent\textbf{\textcolor{biru}{Langkah #1.}}\ }
\newcommand{\jadi}{\par\smallskip\noindent$\Rightarrow$\ \textbf{Jadi, }}
\newcommand{\benar}{\textbf{\textcolor{hijau}{BENAR}}}
\newcommand{\salah}{\textbf{\textcolor{merah}{SALAH}}}
\newcommand{\pernyataan}[1]{\par\smallskip\noindent\textbf{Pernyataan (#1).}\ }
\newcommand{\bagianjudul}[1]{\par\bigskip\noindent%
  {\large\bfseries\color{ungu}\rule[-2pt]{4pt}{16pt}\ #1}\par\smallskip}
\newcommand{\nomorjudul}[2]{\par\bigskip\noindent%
  {\large\bfseries\color{ungu}Pembahasan Nomor #1}\ {\normalfont\itshape\color{gray}(#2)}\par\nopagebreak\smallskip}

\linespread{1.04}

% ============================================================
\begin{document}

% ---------- HALAMAN JUDUL ----------
\thispagestyle{empty}
\begin{center}
\vspace*{1.5cm}
{\color{ungu}\rule{\textwidth}{2pt}}\\[0.5cm]
{\Huge\bfseries\color{ungu} PEMBAHASAN LENGKAP}\\[0.3cm]
{\LARGE\bfseries KUIS 3: MATEMATIKA (MAE)}\\[0.4cm]
{\large\itshape National Card Strategy Championship: Panggung Sang Analis}\\[0.3cm]
{\color{ungu}\rule{\textwidth}{2pt}}\\[1cm]

\begin{tcolorbox}[enhanced,width=0.9\textwidth,colback=ungumuda,colframe=ungu,arc=4pt,boxrule=1pt]
\centering\small
Buku ini membahas \textbf{seluruh 12 soal} secara runtut: konsep kunci, langkah
demi langkah, dan jebakan. Materi: \textbf{Statistika} (Bab 6: distribusi
frekuensi, pemusatan, letak, penyebaran, diagram pencar) dan \textbf{Peluang}
(Bab 7: pencacahan, ruang sampel, peluang kejadian, frekuensi harapan, kejadian
majemuk). Data terbuka bertahap lewat \emph{Berkas A--D}; seluruh angka sinkron
dengan rekap turnamen (rata-rata skor $25$, set remi $52$ kartu).
\end{tcolorbox}

\vfill
{\large Tahun Pelajaran 2026/2027}
\end{center}
\clearpage

\begin{tips}
\textbf{Ringkasan data (dipakai berulang).}
\textbf{Berkas A} (skor $40$ peserta; kelas $0$--$9,\dots,40$--$49$; $f=4,10,12,8,6$;
titik tengah $x=4{,}5;14{,}5;24{,}5;34{,}5;44{,}5$): $\bar{x}=25$; $Me=24{,}5$;
$Mo\approx22{,}83$; $Q_1=15{,}5$; $Q_3=34{,}5$; JAK$=19$; ragam $=144{,}75$;
$s\approx12{,}03$ ($s/\bar{x}\approx48\%$ $\Rightarrow$ beragam/timpang).
\textbf{Berkas B} ($15$ finalis):
korelasi \emph{negatif kuat} ($r\approx-0{,}90$), pencilan \textbf{Bima}; tanpa
Bima menguat ke $r\approx-0{,}99$. \textbf{Berkas C} (remi $52$):
$P(\mathrm{As})=\tfrac{1}{13}$, $P(\text{merah})=\tfrac12$,
$P(\text{gambar})=\tfrac{3}{13}$; frekuensi harapan (dari $1.300$): As $=100$,
merah $=650$, gambar $=300$; tangan $5$ hati $=\tfrac{1287}{2598960}$.
\textbf{Berkas D} ($2,4,8,14,19,23,24,31,38,44$): $\bar{x}=20{,}7$; $Me=21$;
$Q_1=8$, $Q_3=31$.
\end{tips}

% ################################################################
\bagianjudul{BAGIAN A: Memeriksa Hasil Pertandingan (PG Kompleks)}
% ################################################################

\begin{ingat}{Rumus statistika data berkelompok}
Titik tengah $x_i=\tfrac{\text{batas bawah}+\text{batas atas}}{2}$;\quad
$\bar{x}=\dfrac{\sum f_ix_i}{\sum f_i}$;\quad
$Mo=L+\left(\dfrac{d_1}{d_1+d_2}\right)c$;\quad
$Q_i=L+\left(\dfrac{\frac{i\,n}{4}-F}{f}\right)c$ (median $=Q_2$);\quad
ragam $s^2=\dfrac{\sum f_i(x_i-\bar{x})^2}{n}$, simpangan baku $s=\sqrt{s^2}$.
Di sini $L=$ tepi bawah kelas, $F=$ frekuensi kumulatif sebelum kelas, $c=$
panjang kelas.
\end{ingat}

% ================= NOMOR 1 =================
\nomorjudul{1}{6.1 Distribusi Frekuensi: membaca Berkas A}
\begin{soalbox}{Soal 1}
Benar/Salah: (1) titik tengah $10$--$19$ $=14{,}5$; (2) panjang kelas $=10$;
(3) frekuensi kumulatif $<29{,}5$ adalah $26$; (4) kelas modus $=30$--$39$.
\end{soalbox}

\jawab
Tepi kelas: $-0{,}5$; $9{,}5$; $19{,}5$; $29{,}5$; $39{,}5$; $49{,}5$. Frekuensi
kumulatif ($\le$ tepi atas): $4,14,26,34,40$.

\pernyataan{1} Titik tengah $10$--$19=\tfrac{10+19}{2}=14{,}5$. \benar.

\pernyataan{2} Panjang kelas $=9{,}5-(-0{,}5)=10$. \benar.

\pernyataan{3} Frekuensi kumulatif $<29{,}5=4+10+12=26$. \benar.

\pernyataan{4} Frekuensi tertinggi $f=12$ ada pada kelas $20$--$29$, bukan
$30$--$39$ ($f=8$). \salah.

\jadi \textbf{(1) B, (2) B, (3) B, (4) S}.

\begin{jebakan}
Menganggap ``kelas modus'' sekadar kelas terakhir atau kelas dengan angka
terbesar. Kelas modus $=$ kelas dengan \emph{frekuensi} terbesar.
\end{jebakan}

% ================= NOMOR 2 =================
\nomorjudul{2}{6.2 Ukuran Pemusatan: mean, median, modus}
\begin{soalbox}{Soal 2}
Benar/Salah: (1) $\bar{x}=25$ \& $Me=24{,}5$; (2) $Mo=24{,}5$; (3) mean $>$ median
$>$ modus $\Rightarrow$ condong kanan; (4) \emph{klaim} ``rata-rata sudah bagus''
dapat diterima.
\end{soalbox}

\jawab
\pernyataan{1} $\sum fx=4(4{,}5)+10(14{,}5)+12(24{,}5)+8(34{,}5)+6(44{,}5)=1000\Rightarrow\bar{x}=25$;
median $=19{,}5+\tfrac{20-14}{12}\cdot10=24{,}5$. Keduanya benar. \benar.

\pernyataan{2} Kelas modus $20$--$29$: $d_1=12-10=2$, $d_2=12-8=4$,
$Mo=19{,}5+\left(\tfrac{2}{6}\right)10=22{,}83\neq24{,}5$ (nilai $24{,}5$ itu median). \salah.

\pernyataan{3} $\bar{x}=25>Me=24{,}5>Mo=22{,}83\Rightarrow$ sebaran \emph{condong ke
kanan (positif)}. \benar.

\pernyataan{4} \emph{Menilai klaim:} poin minus (makin kecil makin baik) rata-rata
$=25$ tepat di pertengahan skala $0$--$49$; jadi ``sudah bagus'' \emph{tidak}
didukung, hasil hanya biasa-biasa saja. Klaim \textbf{tak} dapat diterima. \salah.

\jadi \textbf{(1) B, (2) S, (3) B, (4) S}.

% ================= NOMOR 3 =================
\nomorjudul{3}{6.3 Ukuran Letak: kuartil}
\begin{soalbox}{Soal 3}
Benar/Salah: (1) $Q_1=15{,}5$; (2) $Q_3=39{,}5$; (3) $Q_2=24{,}5$;
(4) \emph{terapkan ambang}: poin $12\to$ unggulan, poin $40\to$ zona evaluasi.
\end{soalbox}

\jawab
\pernyataan{1} $\tfrac{n}{4}=10$; kelas $Q_1$ $=10$--$19$ ($F=4,f=10$):
$Q_1=9{,}5+\tfrac{10-4}{10}\cdot10=15{,}5$. \benar.

\pernyataan{2} $\tfrac{3n}{4}=30$; kelas $Q_3$ $=30$--$39$ ($F=26,f=8$):
$Q_3=29{,}5+\tfrac{30-26}{8}\cdot10=34{,}5$, \emph{bukan} $39{,}5$. \salah.

\pernyataan{3} $Q_2=Me=24{,}5$ (Nomor 2). \benar.

\pernyataan{4} Unggulan $=$ poin $<Q_1=15{,}5$: $12<15{,}5$ \checkmark. Zona evaluasi
$=$ poin $>Q_3=34{,}5$: $40>34{,}5$ \checkmark. Kedua penerapan benar. \benar.

\jadi \textbf{(1) B, (2) S, (3) B, (4) B}.

\begin{jebakan}
Salah tepi bawah kelas: untuk kelas $30$--$39$, $L=29{,}5$ (bukan $30$). Salah
$L$ menggeser seluruh kuartil.
\end{jebakan}

% ================= NOMOR 4 =================
\nomorjudul{4}{6.4 Ukuran Penyebaran: ragam \& simpangan baku}
\begin{soalbox}{Soal 4}
Benar/Salah: (1) ragam $=144{,}75$; (2) $s\approx12{,}03$; (3) $s>$ JAK;
(4) \emph{menilai persaingan}: $s$ besar dibanding $\bar{x}$ $\Rightarrow$ timpang.
\end{soalbox}

\jawab
Dengan $\bar{x}=25$, hitung $\sum f(x-\bar{x})^2$:
\begin{center}\small
\renewcommand{\arraystretch}{1.15}
\begin{tabular}{@{}c c c c c@{}}
\toprule
$x$ & $f$ & $x-25$ & $(x-25)^2$ & $f(x-25)^2$\\
\midrule
$4{,}5$  & $4$  & $-20{,}5$ & $420{,}25$ & $1681{,}0$\\
$14{,}5$ & $10$ & $-10{,}5$ & $110{,}25$ & $1102{,}5$\\
$24{,}5$ & $12$ & $-0{,}5$  & $0{,}25$   & $3{,}0$\\
$34{,}5$ & $8$  & $9{,}5$   & $90{,}25$  & $722{,}0$\\
$44{,}5$ & $6$  & $19{,}5$  & $380{,}25$ & $2281{,}5$\\
\midrule
\multicolumn{4}{@{}r}{$\sum=$} & $\mathbf{5790}$\\
\bottomrule
\end{tabular}
\end{center}

\pernyataan{1} Ragam $s^2=\tfrac{5790}{40}=144{,}75$. \benar.

\pernyataan{2} $s=\sqrt{144{,}75}\approx12{,}03$. \benar.

\pernyataan{3} $s\approx12{,}03$ \emph{lebih kecil} daripada JAK $=19$; jadi ``lebih
besar'' \salah.

\pernyataan{4} \emph{Menilai persaingan:} $s\approx12$ terhadap $\bar{x}=25$ berarti
simpangan $\approx48\%$ dari rata-rata, \emph{cukup besar}. Banyak skor jauh dari
rata-rata, jadi persaingan \emph{timpang}, bukan seragam. \benar.

\jadi \textbf{(1) B, (2) B, (3) S, (4) B}.

\begin{tips}
\textbf{Ambang ``besar/kecil''.} Simpangan baku tak berarti tanpa pembanding.
Membandingkannya dengan rata-rata ($\tfrac{s}{\bar{x}}\approx48\%$) memberi ukuran
\emph{relatif} keberagaman, cara sederhana menilai timpang/seragamnya sebaran.
\end{tips}

% ================= NOMOR 5 =================
\nomorjudul{5}{6.5 Diagram Pencar: baca Berkas B}
\begin{soalbox}{Soal 5}
Diagram pencar $15$ finalis (waktu berpikir vs poin minus); Bima $(40,26)$.
Benar/Salah: (1) arah negatif; (2) Bima pencilan; (3) pencilan membuat korelasi
hilang; (4) tanpa Bima korelasi makin kuat.
\end{soalbox}

\begin{ingat}{Membaca diagram pencar: arah, kekuatan, pencilan}
Amati tiga hal: \emph{arah} (positif/negatif), \emph{kekuatan} (titik makin rapat
ke garis $\Rightarrow$ makin kuat), dan \emph{pencilan/outlier} (titik yang jauh
menyimpang dari pola). Satu pencilan dapat \emph{memperlemah} korelasi tetapi
tidak selalu menghapus arah hubungan.
\end{ingat}

\jawab
Mayoritas dari $15$ titik menurun saat $x$ naik; hanya \textbf{Bima} $(40,26)$ yang
melonjak ke atas (berpikir lama tetapi poin minus tinggi).

\pernyataan{1} Kecenderungan umum: $x$ naik, $y$ turun $\Rightarrow$ arah
\emph{negatif}. \benar.

\pernyataan{2} Pada $x=40$ tren titik lain memperkirakan $y\approx8$--$9$, tetapi
Bima $y=26$, menyimpang jauh $\Rightarrow$ \emph{pencilan}. \benar.

\pernyataan{3} Meski Bima memperlemah korelasi (dari $r\approx-0{,}99$ menjadi
$\approx-0{,}90$), arah negatif \emph{tetap kuat}; korelasi tidak hilang. \salah.

\pernyataan{4} Tanpa Bima, titik-titik makin rapat pada garis menurun
($r\approx-0{,}99$) $\Rightarrow$ hubungan makin kuat. \benar.

\jadi \textbf{(1) B, (2) B, (3) S, (4) B}.

\begin{tips}
\textbf{Pengaruh pencilan.} Dengan Bima, $r\approx-0{,}90$; tanpa Bima
$r\approx-0{,}99$. Sebuah pencilan menyeret nilai korelasi mendekati nol; karena
itu analis selalu memeriksa pencilan sebelum menyimpulkan kekuatan hubungan. Bima
layak ditelaah khusus (ditindaklanjuti pada Nomor 10d): mengapa berpikir lama tak
menurunkan poin minusnya?
\end{tips}

% ################################################################
\bagianjudul{BAGIAN B: Menguji Sistem Pembagian Kartu (Uraian)}
% ################################################################

\begin{ingat}{Pencacahan \& peluang}
Aturan perkalian: $n$ tahap dengan $n_1,n_2,\dots$ cara $\Rightarrow n_1\cdot
n_2\cdots$. Permutasi $\,_nP_r=\dfrac{n!}{(n-r)!}$; kombinasi
$\,_nC_r=\dfrac{n!}{r!(n-r)!}=\dfrac{{}_nP_r}{r!}$.
Peluang $P(A)=\dfrac{n(A)}{n(S)}$; frekuensi harapan $=N\cdot P(A)$.
\end{ingat}

% ================= NOMOR 6 =================
\nomorjudul{6}{7.1 Aturan Pencacahan: perkalian, permutasi, kombinasi}
\begin{soalbox}{Soal 6}
(a) susunan berurutan $3$ kartu; (b) cara memilih $3$ kartu (tanpa urutan);
(c) tangan $5$ kartu; (d) tangan $5$ hati.
\end{soalbox}

\begin{ingat}{Permutasi atau kombinasi? Uji ``urutan penting?''}
Tanyakan: \emph{jika urutan objek ditukar, apakah hasilnya dianggap berbeda?}
\textbf{Ya} $\Rightarrow$ \textbf{permutasi} $_nP_r=\dfrac{n!}{(n-r)!}$ (mis.\ kartu
dibuka satu per satu, peringkat juara, kode PIN). \textbf{Tidak} $\Rightarrow$
\textbf{kombinasi} $_nC_r=\dfrac{n!}{r!(n-r)!}=\dfrac{_nP_r}{r!}$ (mis.\ tangan
kartu, memilih anggota tim). Keduanya berkait: $_nP_r=r!\times{}_nC_r$, sebab tiap
kombinasi $r$ objek dapat disusun dalam $r!$ urutan.
\end{ingat}

\jawab
\langkah{a} Kartu dibuka \emph{satu per satu dan urutannya dicatat}, jadi susunan
$(A,B,C)$ dianggap \textbf{berbeda} dari $(B,A,C)$ $\Rightarrow$ urutan penting
(\emph{permutasi}). Aturan perkalian: kartu ke-$1$ ada $52$ pilihan; satu terpakai,
ke-$2$ tinggal $51$; ke-$3$ tinggal $50$:
\[
52\times51\times50=132.600={}_{52}P_{3}.
\]

\langkah{b} Kini urutan \emph{tidak} diperhatikan --- memilih $\{A,B,C\}$ sama saja
apa pun urutan pengambilannya ($\Rightarrow$ \emph{kombinasi}). Kuncinya: pada (a),
setiap kelompok $3$ kartu yang sama sudah terhitung $3!=6$ kali (sebagai $6$ urutan
berbeda). Maka bagi hasil (a) dengan $6$:
\[
{}_{52}C_{3}=\frac{{}_{52}P_{3}}{3!}=\frac{132.600}{6}=\mathbf{22.100}.
\]
Hubungan dengan (a): ${}_{52}P_3=3!\times{}_{52}C_3=6\times22.100$.

\langkah{c} \emph{Tangan $5$ kartu} $=$ kumpulan kartu yang dipegang; menerima kartu
dalam urutan berbeda tetap menghasilkan tangan yang \textbf{sama} $\Rightarrow$
kombinasi:
\[
{}_{52}C_{5}=\frac{52\cdot51\cdot50\cdot49\cdot48}{5!}=\mathbf{2.598.960}.
\]

\langkah{d} ``Seluruhnya hati'': pilih $5$ kartu dari \emph{$13$ kartu hati saja}
(urutan tetap tak penting): ${}_{13}C_{5}=\dfrac{13\cdot12\cdot11\cdot10\cdot9}{5!}
=\mathbf{1.287}$. (Peluang tangan ini dihitung di Nomor 7(e).)

\begin{jebakan}
\textbf{Cara membedakan:} tukar urutan objeknya. Kalau hasil \emph{dianggap berbeda}
(kartu dibuka berurutan, peringkat, kode) $\to$ \emph{permutasi}; kalau \emph{tetap
sama} (tangan kartu, kelompok/tim) $\to$ \emph{kombinasi}. Memakai permutasi untuk
``tangan'' $5$ kartu membuat jawaban $5!=120$ kali terlalu besar.
\end{jebakan}

% ================= NOMOR 7 =================
\nomorjudul{7}{7.2--7.3 Ruang Sampel \& Peluang Kejadian}
\begin{soalbox}{Soal 7}
(a) $n(S)$; (b) $P(\mathrm{As})$; (c) $P(\text{merah})$; (d) $P(\text{gambar})$;
(e) peluang tangan $5$ hati (pakai Nomor 6).
\end{soalbox}

\jawab
\langkah{a} Satu kartu dari $52$: $n(S)=\mathbf{52}$.

\langkah{b} $4$ As: $P(\mathrm{As})=\dfrac{4}{52}=\dfrac{1}{13}\approx0{,}077$.

\langkah{c} $26$ merah: $P(\text{merah})=\dfrac{26}{52}=\dfrac{1}{2}$.

\langkah{d} $12$ kartu gambar ($\mathrm{J,Q,K}$):
$P(\text{gambar})=\dfrac{12}{52}=\dfrac{3}{13}\approx0{,}231$.

\langkah{e} Tangan $5$ hati: pakai $_{13}C_5=1.287$ (Nomor 6d) dan
$_{52}C_5=2.598.960$ (Nomor 6c):
\[
P(\text{5 hati})=\frac{{}_{13}C_5}{{}_{52}C_5}=\frac{1.287}{2.598.960}
\approx\mathbf{0{,}000495}\ (\approx0{,}05\%).
\]

\begin{tips}
Peluang kejadian yang ``rumit'' dihitung sama saja: $\dfrac{n(A)}{n(S)}$, dengan
$n(A)$ dan $n(S)$ dicari lewat \emph{kombinatorika} (Nomor 6). Inilah jembatan
antara pencacahan (7.1) dan peluang (7.3).
\end{tips}

% ================= NOMOR 8 =================
\nomorjudul{8}{7.4 Frekuensi Harapan}
\begin{soalbox}{Soal 8}
$N=1.300$ pengambilan (dengan pengembalian). (a) As; (b) merah; (c) gambar;
(d) As sekop; (e) bandingkan dengan pengamatan $92$.
\end{soalbox}

\jawab
Frekuensi harapan $=N\cdot P$.

\langkah{a} As: $1.300\cdot\tfrac{1}{13}=\mathbf{100}$.

\langkah{b} Merah: $1.300\cdot\tfrac12=\mathbf{650}$.

\langkah{c} Gambar: $1.300\cdot\tfrac{3}{13}=\mathbf{300}$.

\langkah{d} As sekop (satu kartu tertentu): $1.300\cdot\tfrac{1}{52}=\mathbf{25}$.

\langkah{e} Selisih $=100-92=\mathbf{8}$, yakni $\dfrac{8}{100}\times100\%=\mathbf{8\%}$
dari frekuensi harapan. Karena $8\%<10\%$ (ambang panitia), tergolong
\textbf{selisih kecil} $\Rightarrow$ frekuensi nyata konsisten dengan model peluang;
\textbf{tudingan ``As lebih jarang keluar'' tidak didukung data}.

\begin{tips}
Makin banyak percobaan $N$, frekuensi relatif $\tfrac{92}{\cdots}$ cenderung makin
dekat ke peluang teoretis $\tfrac{1}{13}$ (\emph{Hukum Bilangan Besar}).
\end{tips}

% ================= NOMOR 9 =================
\nomorjudul{9}{7.5 Kejadian Majemuk}
\begin{soalbox}{Soal 9}
(a) bukan As; (b) As atau King; (c) As atau merah; (d) dua merah dengan
pengembalian; (e) dua merah tanpa pengembalian, bandingkan.
\end{soalbox}

\begin{ingat}{Aturan kejadian majemuk}
Komplemen $P(A^c)=1-P(A)$. Saling lepas: $P(A\cup B)=P(A)+P(B)$. Tidak saling
lepas: $P(A\cup B)=P(A)+P(B)-P(A\cap B)$. Dua kejadian bebas:
$P(A\cap B)=P(A)P(B)$. Tak bebas (bersyarat): $P(A\cap B)=P(A)\,P(B\mid A)$.
\end{ingat}

\jawab
\langkah{a} $P(\text{bukan As})=1-\dfrac{4}{52}=\dfrac{48}{52}=\dfrac{12}{13}$.

\langkah{b} As dan King \emph{saling lepas} (tak ada irisan):
$P(\mathrm{As}\cup\mathrm{K})=\dfrac{4}{52}+\dfrac{4}{52}=\dfrac{8}{52}=\dfrac{2}{13}$.

\langkah{c} As dan merah \emph{tidak saling lepas} (ada $2$ As merah):
\[
P(\mathrm{As}\cup\text{merah})=\dfrac{4}{52}+\dfrac{26}{52}-\dfrac{2}{52}
=\dfrac{28}{52}=\dfrac{7}{13}.
\]

\langkah{d} \emph{Dengan pengembalian:} kartu pertama dikembalikan lalu dikocok,
keadaan kembali seperti semula. Peluang kartu ke-$2$ merah \textbf{tidak bergantung}
pada kartu ke-$1$ $\Rightarrow$ dua kejadian \emph{bebas}, keduanya $\tfrac{26}{52}=\tfrac12$:
\[
P(\text{merah, merah})=\dfrac{26}{52}\cdot\dfrac{26}{52}=\dfrac12\cdot\dfrac12=\dfrac14=0{,}25.
\]

\langkah{e} \emph{Tanpa pengembalian:} kartu pertama \textbf{tidak} dikembalikan,
jadi tersisa $51$ kartu (bila yang pertama merah, sisa merah $=25$). Kejadian ke-$2$
\emph{bergantung} pada ke-$1$ (\emph{tak bebas}/bersyarat),
$P(\text{merah}_2\mid\text{merah}_1)=\tfrac{25}{51}$:
\[
P=\dfrac{26}{52}\cdot\dfrac{25}{51}=\dfrac12\cdot\dfrac{25}{51}=\dfrac{25}{102}\approx0{,}245.
\]
Nilainya \emph{sedikit lebih kecil} daripada (d) ($0{,}245<0{,}25$) karena stok kartu
merah berkurang setelah pengambilan pertama.

\begin{jebakan}
Menganggap ``As atau merah'' saling lepas lalu menjumlah $\tfrac{4}{52}+\tfrac{26}{52}
=\tfrac{30}{52}$. Salah: irisan (As merah) terhitung dua kali, harus dikurangi.
\end{jebakan}

% ################################################################
\bagianjudul{BAGIAN C: Vonis, Apakah Protes Berdasar?}
% ################################################################

% ================= NOMOR 10 =================
\nomorjudul{10}{Sintesis: sebaran $\to$ kuartil-sebagai-peluang $\to$ pencilan}
\begin{soalbox}{Soal 10}
Berkas A ($n=40$), Berkas B (Nomor 5), Berkas C. (a) $P(20$--$29)$, $P(<20)$;
(b) mengapa $P(\text{unggulan})=25\%$ \& mengapa kuartil bukan rata-rata;
(c) frekuensi harapan unggulan, dua peserta tanpa pengembalian, gabungan bebas
dengan As; (d) tindak lanjut pencilan Bima; (e) vonis atas protes.
\end{soalbox}

\begin{ingat}{Ukuran letak \emph{adalah} peluang}
Menurut definisi, kuartil/persentil ke-$p$ adalah nilai yang membuat
$P(X<\text{nilai})=p$. Jadi $Q_1$ memenuhi $P(X<Q_1)=25\%$: sebuah \emph{ukuran
letak} (Bab 6) sekaligus sebuah \emph{peluang} (Bab 7).
\end{ingat}

\jawab
Frekuensi kelas: $0$--$9\,(4)$, $10$--$19\,(10)$, $20$--$29\,(12)$, $30$--$39\,(8)$,
$40$--$49\,(6)$; $n=40$.

\langkah{a} Peluang empiris (frekuensi relatif):
\[
P(20\text{--}29)=\frac{12}{40}=0{,}3;\qquad
P(\text{skor}<20)=\frac{4+10}{40}=\frac{7}{20}=0{,}35.
\]

\langkah{b} \textbf{(i)} Gelar unggulan $=$ berada di bawah $Q_1$. Menurut
\emph{definisi} kuartil, tepat $25\%$ data ada di bawah $Q_1$, berapa pun
frekuensi tiap kelasnya. Maka peluang peserta acak bergelar unggulan
$=P(X<Q_1)=\mathbf{25\%}$, \emph{tanpa} menghitung frekuensi lagi.
\textbf{(ii)} Sebaran \emph{condong ke kanan} (Nomor 2: $\bar{x}=25>Me=24{,}5>Mo$),
sehingga rata-rata $25$ \emph{tertarik} oleh sedikit peserta ber-poin minus besar
dan tak lagi mewakili ``$25\%$ terbaik''. Kuartil bersifat \emph{posisional}:
selalu memotong tepat $25\%$ data dan \emph{tahan} terhadap nilai ekstrem; itu
sebabnya kuartil dipakai sebagai ambang, bukan rata-rata.

\langkah{c} \textbf{(i)} Unggulan dari $200$: $200\times\tfrac14=\mathbf{50}$.
\textbf{(ii)} Dua peserta \emph{tanpa} pengembalian:
$P=\dfrac{10}{40}\cdot\dfrac{9}{39}=\dfrac14\cdot\dfrac{3}{13}=\dfrac{3}{52}\approx0{,}058$.
\textbf{(iii)} Peserta \& kartu \emph{bebas}, kalikan (dengan $P(\mathrm{As})=\tfrac1{13}$):
$P(\text{unggulan}\cap\mathrm{As})=\dfrac14\cdot\dfrac1{13}=\dfrac{1}{52}\approx0{,}019$.

\langkah{d} \textbf{(i)} \emph{Tidak}. Dari $15$ finalis, $14$ mengikuti pola
``berpikir lebih lama $\to$ poin minus lebih kecil''; korelasi keseluruhan tetap
kuat ($r\approx-0{,}90$). Satu pencilan tak membatalkan kesimpulan
\emph{keterampilan}; Bima hanya kasus khusus untuk ditelaah.
\textbf{(ii)} Mengeluarkan Bima membuat korelasi \textbf{menguat}
($r\approx-0{,}90\to-0{,}99$): titik Bima yang menyimpang tadinya menyeret $r$
mendekati nol; tanpa Bima titik-titik makin rapat ke garis.

\langkah{e} \emph{Vonis atas protes} (lewat dua indikator panitia).
\emph{Indikator I, keseimbangan peluang:} As nyata $92$ vs harapan $100$
(Nomor 8; selisih $8\%<10\%$), $P(\text{merah})=\tfrac12$ seimbang $\Rightarrow$
\textbf{terpenuhi}. \emph{Indikator II, sebaran terkendali \& keterampilan:}
$s\approx12{,}03$, JAK $=19$ (Nomor 3--4) terkendali; korelasi $r\approx-0{,}90$
(Nomor 5) condong ke keterampilan $\Rightarrow$ \textbf{terpenuhi}.
\jadi \textbf{protes TIDAK berdasar}: tak ada bukti pembagian kartu berat sebelah,
dan hasil yang ``timpang'' justru mencerminkan perbedaan \emph{keterampilan}
(pemain yang berpikir lebih cermat berpoin minus lebih kecil), bukan
ketidakadilan. Bukti: $\ge2$ statistik ($s\approx12{,}03$, JAK $=19$,
$r\approx-0{,}90$) dan $\ge1$ peluang (harapan As $100$ vs nyata $92$).
\textbf{Saran:} perbanyak ronde/pengamatan (Hukum Bilangan Besar), telaah pencilan
Bima, dan umumkan hasil uji keadilan ini agar peserta menerima; musim depan
bebas protes serupa.

\begin{tips}
\textbf{Momen ``menggabung''.} $Q_1$, sebuah \emph{ukuran letak} (Bab 6),
langsung memberi \emph{peluang} $25\%$ (Bab 7): persentil \emph{adalah} peluang.
Statistik menata data; peluang membacanya sebagai kesempatan. Dua bab yang tampak
terpisah ternyata satu bahasa.
\end{tips}

\begin{jebakan}
Pada (c)(ii) memakai peluang \emph{dengan} pengembalian
$\left(\tfrac14\right)^2=\tfrac{1}{16}$. Karena peserta tak dikembalikan,
pengambilan kedua memakai $\tfrac{9}{39}$, sehingga
$\tfrac{3}{52}\approx0{,}058$, sedikit lebih kecil.
\end{jebakan}

% ################################################################
\bagianjudul{BAGIAN D: Soal Bonus}
% ################################################################

% ================= NOMOR 11 =================
\nomorjudul{11}{[BONUS] Statistik ``Sepuluh Kartu Penentu Kemenangan''}
\begin{soalbox}{Soal 11}
Nomor $10$ kartu penentu (dari $52$ kartu bernomor $1$--$52$):
$2,4,8,14,19,23,24,31,38,44$ ($n=10$). (a) mean; (b) median \& modus;
(c) $Q_1,Q_3$, JAK; (d) jangkauan.
\end{soalbox}

\jawab
\langkah{a} $\sum x=2+4+8+14+19+23+24+31+38+44=207$,
maka $\bar{x}=\dfrac{207}{10}=\mathbf{20{,}7}$.

\langkah{b} $n=10$ (genap): median $=\dfrac{x_5+x_6}{2}=\dfrac{19+23}{2}=\mathbf{21}$.
Semua nilai berbeda $\Rightarrow$ \textbf{tidak ada modus}.

\langkah{c} Karena $n=10$ \emph{genap}, median ($21$) jatuh di antara dua data,
sehingga data langsung terbagi rata menjadi $5$ nilai bawah dan $5$ nilai atas.
$Q_1$ = median kelompok bawah, $Q_3$ = median kelompok atas:
\[
\underbrace{2,\;4,\;\mathbf{8},\;14,\;19}_{Q_1=8}\qquad\qquad
\underbrace{23,\;24,\;\mathbf{31},\;38,\;44}_{Q_3=31}.
\]
JAK $=Q_3-Q_1=31-8=\mathbf{23}$.

\langkah{d} Jangkauan $=44-2=\mathbf{42}$. Interpretasi: nomor kartu penentu
\emph{tersebar luas} (dari $2$ hingga $44$ dari kemungkinan $1$--$52$), tidak
terpusat pada satu wilayah; kemenangan tidak bergantung pada segelintir kartu
saja.

% ================= NOMOR 12 =================
\nomorjudul{12}{[BONUS] Peluang Kombinatorik: menguji keadilan tangan}
\begin{soalbox}{Soal 12}
(a) peluang dua kartu berpasangan; (b) peluang tangan $5$ berisi tepat dua As;
(c) interpretasi keadilan.
\end{soalbox}

\jawab
\langkah{a} Kartu pertama bebas; agar berpasangan, kartu kedua harus satu rank
dengan kartu pertama. Sisa kartu serank $=3$ dari $51$:
\[
P(\text{pasangan})=\dfrac{3}{51}=\dfrac{1}{17}\approx0{,}059.
\]

\langkah{b} Tangan $5$ kartu $\Rightarrow$ urutan tak penting (\emph{kombinasi}).
``Tepat dua As'': pilih $2$ dari $4$ As \textbf{dan} $3$ dari $48$ kartu non-As
(aturan perkalian pada kombinasi). Penyebut $\binom{52}{5}$ = seluruh tangan yang
mungkin:
\[
P=\dfrac{\binom{4}{2}\binom{48}{3}}{\binom{52}{5}}
=\dfrac{6\times17.296}{2.598.960}
=\dfrac{103.776}{2.598.960}\approx\mathbf{0{,}0399}\ (\approx3{,}99\%).
\]

\langkah{c} \emph{Interpretasi.} Peluang $\approx4\%$ tergolong \textbf{kecil},
sehingga memperoleh ``tepat dua As'' \emph{jarang} terjadi. Karena peluang ini
sama bagi setiap pemain dan tidak besar, tak ada pemain yang mudah memperoleh
tangan istimewa $\Rightarrow$ permainan tetap \textbf{seimbang dan adil}.

\begin{tips}
Bonus ini menautkan \emph{kombinatorika} (Bab 7.1) dengan \emph{peluang}
(Bab 7.3): menghitung peluang kejadian rumit $=$ membandingkan banyak
\emph{kombinasi} kejadian yang diinginkan terhadap seluruh kemungkinan.
\end{tips}

\vfill
\begin{center}\small\itshape
Selesai. Semoga bermanfaat sebagai bahan belajar.
\end{center}

\end{document}
